% GENERATED FILE. Edit canonical lessons, then run npm run generate:books.
% canonical-source-hash: fnv1a64:ca2cfd6ab4b80ab1
% canonical-lessons: GU-C150-javano, GU-C150-karvana-cho, GU-C150-kadach, GU-C150-plans-practice, GU-R150-not-ask-past

\chapter{Plans and perhaps}
\label{ch:gu-plans-and-perhaps}
\hlchaptermodality{150}

\begin{chapteropening}
\textbf{By the end of this chapter:} \emph{I can say what I am going to do as a plan already made, ask other people about their plans, call a plan off, and say what will perhaps happen.}

Carry a short exchange about tomorrow that holds a plan, a plan called off, a guess and a promise.
\end{chapteropening}

\section[\texorpdfstring{…\gu{વાનો} \gu{છું} (-vāno chũ)}{…વાનો છું (-vāno chũ)}]{\gu{હું} \gu{જવાનો} \gu{છું} — a plan already made}

\label{lesson:GU-C150-javano}



\begin{quote}
Say \emph{I'll go to the market tomorrow}. Then name the ending every verb's name wears: \textbf{-\gu{વું}}, as in \textbf{\gu{જવું}}.
\end{quote}

\begin{grammarlens}[title={Grammar Lens: -\gu{વાનો} and the copula}]
\textbf{\gu{હું} \gu{જઈશ}} says \emph{I will go}. \textbf{\gu{હું} \gu{જવાનો} \gu{છું}} says \emph{I'm going to go}: the plan is already made.

\begin{quote}
\textbf{\gu{હું} \gu{કાલે} \gu{જવાનો} \gu{છું}.} — \emph{hũ kāle javāno chũ.} — I'm going tomorrow. (a man) \textbf{\gu{હું} \gu{કાલે} \gu{જવાની} \gu{છું}.} — \emph{hũ kāle javānī chũ.} — the same, said by a woman.
\end{quote}

The pieces: the verb's name \textbf{\gu{જવું}}, with \textbf{-\gu{વું}} turned into \textbf{-\gu{વાનો}}, then the present copula, \textbf{\gu{છું}}, \textbf{\gu{છે}}, \textbf{\gu{છીએ}} or \textbf{\gu{છો}}.

\textbf{-\gu{વાનો}} agrees like \emph{my}: \textbf{-\gu{વાનો}}, \textbf{-\gu{વાની}}, \textbf{-\gu{વાનું}}, \textbf{-\gu{વાના}}. So a group says \textbf{\gu{અમે} \gu{જવાના} \gu{છીએ}}, \emph{we're going to go}.
\end{grammarlens}

\subsection*{Your turn}
\begin{itemize}
\raggedright
  \item \emph{Say it:} \emph{hũ kāle javāno chũ}, or \emph{javānī}
  \item \emph{Say it:} \textbf{\gu{તે} \gu{કામ} \gu{કરવાનો} \gu{છે}} — \emph{te kām karvāno che}, he is going to work
  \item \emph{Contrast:} \emph{hũ jaīś} and \emph{hũ javāno chũ}: which one is a plan already made
\end{itemize}

\begin{tcolorbox}[breakable,colback=teal!4,colframe=teal!35!black,title={Before you move on}]
Which part of \textbf{\gu{જવાનો}} agrees with the speaker? (The last sign: \textbf{-\gu{ો}} or \textbf{-\gu{ી}}.) Which copula follows it? (The present one: \textbf{\gu{છું}}.)
\end{tcolorbox}
\section[\texorpdfstring{\gu{શું} \gu{કરવાના} \gu{છો}? (śũ karvānā cho?)}{શું કરવાના છો? (śũ karvānā cho?)}]{\gu{તમે} \gu{શું} \gu{કરવાના} \gu{છો}? — asking about plans}

\label{lesson:GU-C150-karvana-cho}



\begin{quote}
Say \emph{I'm going tomorrow}, as a plan already made. Then ask \emph{do you speak Gujarati?} with \textbf{\gu{શું}} in front.
\end{quote}

\begin{grammarlens}[title={Grammar Lens: the question}]
\begin{quote}
\textbf{\gu{તમે} \gu{આવતીકાલે} \gu{શું} \gu{કરવાના} \gu{છો}?} — \emph{tame āvtīkāle śũ karvānā cho?} — What are you going to do tomorrow?
\end{quote}

\textbf{\gu{તમે}} takes \textbf{-\gu{વાના}} and \textbf{\gu{છો}}: respect takes the \textbf{-\gu{ા}} forms here as it does with \textbf{\gu{મારા} \gu{પિતા}}. Put \textbf{\gu{શું}} at the front instead of inside, and the plan becomes a yes-or-no question:

\begin{quote}
\textbf{\gu{શું} \gu{તમે} \gu{આવવાના} \gu{છો}?} — \emph{śũ tame āvvānā cho?} — Are you going to come?
\end{quote}
\end{grammarlens}

\subsection*{You'll want to know: the plan is off}
\textbf{\gu{નથી}} takes the copula's place, as in every present negative:

\begin{quote}
\textbf{\gu{ના}, \gu{હું} \gu{નથી} \gu{આવવાનો}.} — \emph{nā, hũ nathī āvvāno.} — No, I'm not going to come.
\end{quote}

\subsection*{Your turn}
\begin{itemize}
\raggedright
  \item \emph{Say it:} \emph{tame śũ karvānā cho?}
  \item \emph{Answer:} with \textbf{\gu{હું} … \gu{જવાનો} \gu{છું}} or \textbf{\gu{જવાની} \gu{છું}}
  \item \emph{Say it:} \emph{hũ nathī āvvāno}, or \emph{āvvānī}
\end{itemize}

\begin{tcolorbox}[breakable,colback=teal!4,colframe=teal!35!black,title={Before you move on}]
What does \textbf{\gu{તમે}} take before \textbf{\gu{છો}}? (\textbf{-\gu{વાના}}.) Where does \textbf{\gu{નથી}} go when a plan is off? (In the copula's place: \textbf{\gu{હું} \gu{નથી} \gu{આવવાનો}}.)
\end{tcolorbox}
\section[\texorpdfstring{\gu{કદાચ} (kadāc)}{કદાચ (kadāc)}]{\gu{કદાચ} (kadāc) — perhaps}

\label{lesson:GU-C150-kadach}



\begin{quote}
Before the new one: ask \emph{what are you going to do tomorrow?}, then say \emph{I'm not going to come}.
\end{quote}

\subsection*{You'll want to know: \gu{કદાચ}}
\textbf{\gu{કદાચ}} (\emph{kadāc}) — \textquotedblleft{}perhaps\textquotedblright{}.

A plan is sure; a guess is not. \textbf{\gu{કદાચ}} usually stands at the front, and it pairs easily with the future:

\begin{quote}
\textbf{\gu{કદાચ} \gu{તે} \gu{આવશે}.} — \emph{kadāc te āvśe.} — Perhaps she'll come. \textbf{\gu{કદાચ} \gu{આજે} \gu{વરસાદ} \gu{આવશે}.} — \emph{kadāc āje varsād āvśe.} — Maybe it will rain today.
\end{quote}

\begin{cousinweb}
\textbf{\gu{કદાચ}} (\emph{kadāc}) is usually traced to Sanskrit \emph{kadācit}, \textquotedblleft{}at some time, once\textquotedblright{}, built on \emph{kadā}, \textquotedblleft{}when\textquotedblright{}. Sanskrit already let \emph{kadācit} lean towards \textquotedblleft{}perhaps\textquotedblright{}, and that is the sense Gujarati kept.
\end{cousinweb}

\begin{grammarlens}[title={What you've built}]
One more word: the future can now be sure, planned, or only possible.
\end{grammarlens}

\subsection*{Your turn}
Say these aloud:

\begin{itemize}
\raggedright
  \item \emph{kadāc}
  \item \emph{kadāc te āvśe}
  \item \textbf{\gu{કદાચ} \gu{હું} \gu{મોડો} \gu{આવીશ}} — \emph{kadāc hũ moḍo āvīś}, perhaps I'll come late
\end{itemize}

\begin{tcolorbox}[breakable,colback=teal!4,colframe=teal!35!black,title={Before you move on}]
What does \textbf{\gu{કદાચ}} mean? (\emph{Perhaps}.) Which tense does it pair with most easily? (The future.)
\end{tcolorbox}
\section[Practice]{Sure, planned, possible — a practice}

\label{lesson:GU-C150-plans-practice}



\begin{quote}
Say \emph{perhaps she'll come}.
\end{quote}

\subsection*{The exchange: sure, planned, possible}
\begin{quote}
\textbf{— \gu{તમે} \gu{આવતીકાલે} \gu{શું} \gu{કરવાના} \gu{છો}?} \textbf{— \gu{હું} \gu{બજાર} \gu{જવાની} \gu{છું}. \gu{તમે}?} \textbf{— \gu{હું} \gu{ઘરે} \gu{રહેવાનો} \gu{છું}. \gu{હું} \gu{બજાર} \gu{નથી} \gu{જવાનો}.} \textbf{— \gu{કદાચ} \gu{મારી} \gu{બહેન} \gu{આવશે}.} \textbf{— \gu{તો} \gu{સાંજે} \gu{આપણે} \gu{ચા} \gu{પીશું}.}
\end{quote}

Say these aloud:

\begin{itemize}
\raggedright
  \item the exchange aloud, taking each part in turn
  \item which line is a plan, which is a guess, and which is a promise
\end{itemize}

\textbf{\gu{જવાની} \gu{છું}} and \textbf{\gu{રહેવાનો} \gu{છું}} are plans already made. \textbf{\gu{કદાચ} … \gu{આવશે}} is a guess. \textbf{\gu{પીશું}} is a plain future, close to a promise.

\subsection*{Your turn}
\begin{itemize}
\raggedright
  \item \emph{Answer:} \textbf{\gu{તમે} \gu{આવતીકાલે} \gu{શું} \gu{કરવાના} \gu{છો}?} with one plan and one guess
  \item \emph{Say it:} the plain-future question beside it: \textbf{\gu{તમે} \gu{ક્યારે} \gu{આવશો}?} — \emph{tame kyāre āvśo?}
\end{itemize}

\begin{tcolorbox}[breakable,colback=teal!4,colframe=teal!35!black,title={Before you move on}]
Which word turns a statement about tomorrow into a guess? (\textbf{\gu{કદાચ}}.) Which form says the plan is already made? (\textbf{-\gu{વાનો} \gu{છું}}, or \textbf{-\gu{વાની} \gu{છું}}.)
\end{tcolorbox}
\section[(dialogue)]{Not, yes or no, and the past — a review}

\label{lesson:GU-R150-not-ask-past}



\begin{quote}
Say \emph{perhaps she'll come}, then \emph{I'm going tomorrow}, as a plan, then \emph{I'm sleepy}.

\begin{itemize}
\raggedright
  \item \emph{Write it:} \textbf{\gu{ઓ}}, \textbf{\gu{ડ}} and \textbf{\gu{ઊ}}, once each, from memory
\end{itemize}
\end{quote}

\subsection*{Your turn: not, and yes or no}
Say these aloud:

\begin{itemize}
\raggedright
  \item \emph{they}; \emph{she speaks Gujarati}; \emph{we speak Gujarati}; \emph{you speak Gujarati}
  \item \emph{I don't speak Gujarati}; \emph{don't worry}; \emph{not now}
  \item \emph{nobody is here}; \emph{I'm not eating anything}; \emph{she never drinks tea}
  \item \emph{do you speak Gujarati?}, with and without \textbf{\gu{શું}}; \emph{are you coming or not?}; \emph{you live here, right?}
  \item \emph{I have a fever}; \emph{they like Gujarati}; \emph{my sister}; \emph{her name is Mira}; \emph{our village}, yours and mine
\end{itemize}

\subsection*{Your turn: the past}
Say these aloud:

\begin{itemize}
\raggedright
  \item \emph{I was there}; \emph{the house was small}; \emph{yesterday, at home}
  \item \emph{I went}; \emph{she came}; \emph{I got worried}; \emph{I came late}
  \item \emph{I drank tea}; \emph{she drank water}; \emph{they drank milk}
  \item \emph{I did the work}; \emph{I ate rice}; \emph{I took the medicine}
\end{itemize}

\begin{tcolorbox}[breakable,colback=teal!4,colframe=teal!35!black,title={Before you move on}]
\emph{They}? (\textbf{\gu{તેઓ}}.) \emph{I don't speak Gujarati}? (\textbf{\gu{હું} \gu{ગુજરાતી} \gu{નથી} \gu{બોલતો}}, or \textbf{\gu{બોલતી}}.) \emph{Do you speak Gujarati?} (\textbf{\gu{શું} \gu{તમે} \gu{ગુજરાતી} \gu{બોલો} \gu{છો}?}) \emph{I have a fever}? (\textbf{\gu{મને} \gu{તાવ} \gu{છે}}.) \emph{I went}? (\textbf{\gu{હું} \gu{ગયો}}, or \textbf{\gu{ગઈ}}.) \emph{Late}? (\textbf{\gu{મોડું}}.) \emph{I drank tea}? (\textbf{\gu{મેં} \gu{ચા} \gu{પીધી}}.) \emph{I ate rice}? (\textbf{\gu{મેં} \gu{ભાત} \gu{ખાધો}}.)
\end{tcolorbox}
